% Member 3 - final report, IV.G What the Experiments Do Not Show
\subsection{What the Experiments Do Not Show}
\label{sec:negative_results}

A study that reports only the comparisons it expected to succeed is hard
to evaluate. This section collects the results that did \emph{not} go our
way, together with the ones that were inconclusive, and states what they
do and do not license.

\paragraph{The spatial modulation does not improve the accuracy.}
This is the central negative result, and it is stated first because it is
the one the reader is most likely to expect to be the opposite. The
strength-matched uniform control and the spatial variant are statistically
indistinguishable (Section~\ref{sec:main_results}), and the equal-strength
comparison agrees. Once the amount of style is held fixed, the spatial
distribution of the strength contributes nothing that five seeds can
detect. We report the paired differences, their confidence intervals and
their $p$ values; the intervals are narrow enough (about $\pm1.2$ points)
that the absence of an effect is informative rather than a lack of power.

\paragraph{The mechanism does not work as designed.}
The second negative result is from the image analysis
(Section~\ref{sec:mechanism}). The modulation was meant to preserve
structure where the structure matters, and it does not: the images with
the most edges do not gain more edge retention than the images with the
fewest ($t = +0.84$, $p = 0.40$), and the per-class difference is
remarkably uniform. The modulation behaves like a mild, nearly uniform
reduction of the style strength. That is a genuine failure of the stated
mechanism, and it is consistent with the classification result: a
mechanism that does not operate cannot produce an effect.

\paragraph{Data volume is not free.}
The +Duplicate control reaches $37.7\%$, which is \emph{below} the
Original group's $40.8\%$: doubling the number of training rows with
byte-identical copies does not merely fail to help, it hurts by about
three points. We read this as the optimiser spending its fixed step budget
on repeated data, which is the expected consequence of a step budget rather
than an epoch budget. It is reported because it is the reason the styled
groups' gain has to be measured against a data-volume control and not
against the Original group alone.

\paragraph{The generic appearance control is inconclusive.}
The +NonStyle group reaches $39.9\%$, and the difference from +Duplicate
($+2.19$ points) does not survive the false-discovery-rate correction
across the family of comparisons. A generic appearance change, as we
implemented it, is therefore not measurably better than a copy. What
\emph{is} significant is the comparison of the styled group against
+NonStyle ($+7.48$ points, $p = 0.0001$), which is the attribution of
style itself rather than of any generic transformation.

\paragraph{What the study does support.}
Three findings survive the controls. Style augmentation itself works: the
styled groups gain about seven points over the data-volume controls, and the
effect is far outside the seed-to-seed spread. The amount of style applied
is the variable that predicts the accuracy, and it does so across the whole
sweep, with the spatial and the uniform configurations lying on the same
curve. And the gain is not a data-volume artefact, because the controls
that hold the data volume fixed do not show it.

\paragraph{Scope.}
These conclusions are about single-source domain generalisation on PACS,
with one source domain, seven classes and a pretrained ResNet-18, at one
operating point of the augmentation strength. We do not claim that the
spatial distribution of the style never matters; we claim that on this
benchmark, at this strength, with this structure map, it does not matter
measurably, and that reporting a gain without controlling for the amount
of style would have made the opposite claim appear to hold.
